\chapter{Diagrama de clases}
\label{cap:clases}

\sloppy

El diagrama de clases (Figura~\ref{fig:clases}) es una propuesta de diseño inferida a partir de la
funcionalidad observada. Se organiza en tres grupos:

\begin{itemize}
    \item \textbf{Datos maestros y seguridad:} \clase{Usuario}, \clase{Rol}, \clase{Permiso},
          \clase{Cliente}, \clase{Trabajador}, \clase{Cuadrilla} y \clase{TipoBulto}. Los permisos
          se asignan a los roles y no al usuario, de modo que \clase{Usuario} solo conoce su rol y
          no concentra la lógica de autorización de cada módulo. Un usuario puede estar ligado a la
          ficha de un \clase{Trabajador} (el jefe tarjador, para conocer sus operaciones asignadas)
          o de un \clase{Cliente} (para consultar solo sus operaciones).
    \item \textbf{Carga y operación en terreno:} \clase{Manifiesto}, \clase{Contenedor},
          \clase{LineaCarga}, \clase{Programacion}, \clase{Operacion}, \clase{SelloSeguridad},
          \clase{Incidencia}, \clase{Fotografia} y \clase{Reporte}.
    \item \textbf{Controladores y servicios:} \clase{GestorManifiesto},
          \clase{GestorProgramacion}, \clase{GestorOperacion} y \clase{GestorReportes} coordinan
          los casos de uso; \clase{ServicioAduana} y \clase{ServicioCorreo} encapsulan los sistemas
          externos; y \clase{SincronizadorOffline} gestiona el envío diferido de fotografías.
\end{itemize}

Las relaciones representan las decisiones de diseño más importantes: un manifiesto se compone de
contenedores y cada contenedor de líneas de carga, cada una de las cuales pertenece a un cliente y
a un tipo de bulto; una programación corresponde a un contenedor, se asigna a una cuadrilla y a un
jefe tarjador (\clase{Trabajador}) y, al iniciarse la faena, da origen a una operación, por lo
que estas asociaciones admiten cero mientras la programación está pendiente; una cuadrilla agrupa a sus
trabajadores; una operación se compone de fotografías, de un sello de seguridad y de las
incidencias registradas, que se respaldan con fotografías de la misma operación (agregación) y
se asocian a la línea de carga afectada, salvo la de sello alterado; el sincronizador mantiene la cola de
fotografías pendientes de envío; y un reporte se genera a partir de una operación y, una vez aprobado, queda asociado al usuario
que lo aprobó. Las dependencias (líneas punteadas) indican qué clases utiliza cada controlador en
los diagramas de secuencia, y \clase{Imagen}, \clase{Filtro} y \clase{File} son tipos de apoyo que
solo aparecen en las firmas de las operaciones. Las operaciones de crear, consultar, editar, eliminar y cambiar estado de los datos
maestros son las habituales de cada entidad y no se modelaron como controladores para no recargar
el diagrama.

\begin{landscape}
\null\vfill
\begin{figure}[H]
    \centering
    \IfFileExists{Figuras/diagramas/21-diagrama-clases.png}{%
        \includegraphics[width=\linewidth]{21-diagrama-clases.png}%
    }{%
        \fbox{\parbox[c][7cm][c]{0.85\linewidth}{\centering
            \textbf{[Diagrama pendiente]}\\[4pt]
            \texttt{Figuras/diagramas/21-diagrama-clases.png}}}%
    }
    \caption[Diagrama de clases]{Diagrama de clases del sistema BeeTracer.}
    \label{fig:clases}
\end{figure}
\vfill
\end{landscape}

\fussy
